\documentclass[aspectratio=169]{beamer}

\usetheme{Madrid}
\usepackage{booktabs}
\usepackage{listings}
\usepackage{xcolor}
\setbeamertemplate{navigation symbols}{}
\setbeamertemplate{caption}{\raggedright\insertcaption\par}

% compact JSON-ish listing style (matches week 4)
\definecolor{kf}{HTML}{1f6feb}
\definecolor{cm}{HTML}{6a737d}
\lstdefinestyle{json}{
  basicstyle=\ttfamily\scriptsize,
  breaklines=true, columns=fullflexible,
  showstringspaces=false,
  keywordstyle=\color{kf}, commentstyle=\color{cm}\itshape,
  morecomment=[l]{//}, morekeywords={},
  frame=none, aboveskip=2pt, belowskip=2pt
}

\title[Model \& dataset card database]{The LULC model \& dataset card database}
\author[Salil Gujar]{Salil Gujar \\ \footnotesize M.Tech CSE \quad Entry No.\ 2025MCS2106}
\date{June 2026}

\begin{document}

\frame{\titlepage}

% ------------------------------------------------------------------
\begin{frame}{Where we are, and what was advised}
\begin{columns}[T]
\begin{column}{0.5\textwidth}
\textbf{Done so far}
\begin{itemize}\small
  \item Weeks 2 and 3: a 4-class base map (Alpha Earth, 10\,m) plus a \textbf{living hierarchy}:
  split / add a class, retrain on the fly.
  \item Week 4: \emph{designed} the \textbf{Model Card} and \textbf{Dataset Card} schema and a
  catalogue, on paper, no code.
\end{itemize}
\end{column}
\begin{column}{0.5\textwidth}
\textbf{What was advised this week}
\begin{itemize}\small
  \item \textbf{Convert the design into code:} a database that stores all the cards.
  \item A dataset needs a \textbf{type}: \emph{training} (just polygons) vs \emph{inference}
  (just the input features).
  \item A model card carries \textbf{both}: its training data and its inference feature source.
  \item Keep \texttt{extent} a simple \textbf{bounding box} for now; keep improving the front end.
  \item Try splitting trees into \textbf{tea / non-tea}.
\end{itemize}
\end{column}
\end{columns}
\end{frame}

% ------------------------------------------------------------------
\begin{frame}{From a design to a running database}
The schema we designed last week is now code the tool writes to on every retrain:
\begin{itemize}
  \item \textbf{Two card schemas:} extended with the new \texttt{type} field and the model's two
  dataset references, plus a validator.
  \item \textbf{A catalogue that \emph{is} the database:} it validates, writes and indexes cards,
  mints them from live training runs, and answers ``which models cover my area''.
  \item The catalogue is a \textbf{git repository}; publishing pushes the cards to GitHub.
  \item New API endpoints and a full-screen \textbf{Model Zoo} browser on the front end.
\end{itemize}
\end{frame}

% ------------------------------------------------------------------
\begin{frame}{The pipeline: retrain $\rightarrow$ card $\rightarrow$ git}
One left-to-right flow, all reusing what we already had:
\begin{enumerate}
  \item \textbf{Retrain a node:} the existing split / add trainer.
  \item \textbf{Mint cards:} lift the held-out report, the classes, and the marked polygons into
  a Model Card and its Dataset Cards.
  \item \textbf{Store them in the catalogue}, validated on write.
  \item \textbf{Publish = git push:} commit the cards and push to the shared GitHub repository.
\end{enumerate}
\end{frame}

% ------------------------------------------------------------------
\begin{frame}{The week-5 change: typed datasets + two references}
\begin{itemize}
  \item \textbf{Training dataset:} the \emph{labels} a model learned from
  (\texttt{type:\,training}). Polygons or pixel tables, with classes; no embedding needed.
  \item \textbf{Inference dataset:} the \emph{input features} a model runs on
  (\texttt{type:\,inference}). The Alpha Earth 64-d embedding; no labels.
  \item A \textbf{Model Card links to both}: \texttt{training.datasets} (what it was taught) and
  \texttt{inference.dataset} (what it consumes), so a model is portable and auditable.
\end{itemize}
\end{frame}

% ------------------------------------------------------------------
\begin{frame}{The database is GitHub}
\begin{itemize}
  \item The catalogue is a \textbf{git working tree} of a shared zoo repository, the same model
  Hugging Face uses for its card hub.
  \item \textbf{Publish = commit + push.} A user trains a model, inspects its card, and pushes it to
  the zoo; versioning, sharing and history come for free.
  \item Only the \textbf{JSON cards} are committed; the trained model binaries stay local. A card
  carries the metrics and provenance that make it useful without the binary.
  \item Live now: cards push to a real GitHub repository.
\end{itemize}
\end{frame}

% ------------------------------------------------------------------
\begin{frame}{Browse, and pick a model for your area}
\begin{columns}[T]
\begin{column}{0.5\textwidth}
\textbf{The catalogue answers queries}
\begin{itemize}\small
  \item List every card in the zoo.
  \item Filter to the models whose extent \textbf{covers your area}.
  \item Open any card in full.
\end{itemize}
\end{column}
\begin{column}{0.5\textwidth}
\textbf{Extent: show the real footprint}
\begin{itemize}\small
  \item For \textbf{polygon data}, the extent \emph{is} the polygons: we draw the prominent ones
  on the map, not a country-sized box (everything here is India, so ``India-wide'' says nothing).
  \item Feature sources (Alpha Earth / Tessera) are genuinely India-wide, so those show a
  \emph{label}. Querying ``is there a model for my area'' is still a plain bbox overlap.
\end{itemize}
\end{column}
\end{columns}
\end{frame}

% ------------------------------------------------------------------
\begin{frame}{Serving the map: Earth Engine tiles, not a PNG}
\begin{itemize}
  \item The base model is linear, so it replays \textbf{as band math inside Earth Engine}, classified
  at native 10\,m on the server.
  \item Realistic mode is served as \textbf{EE map tiles} (an XYZ url), not a fixed PNG: crisp at
  \emph{any} zoom, \emph{any} area size, with nothing downloaded to us.
  \item This is the ``tile URL'' ask: the model is \texttt{expressible\_as\_bandmath}, the hook
  for pushing a model to EE and serving its output as tiles.
\end{itemize}
\end{frame}

% ------------------------------------------------------------------
\begin{frame}{The full-screen Model Zoo}
A dedicated browser, so the cards aren't crammed into the sidebar:
\begin{itemize}
  \item \textbf{Tabs:} Models and Datasets; a responsive grid of uniform card tiles (class
  swatches and a published / local pill).
  \item \textbf{Detail pane:} the classes it produces (with the legend), a per-class metrics table,
  and the linked \textbf{training} and \textbf{inference} datasets as clickable chips.
  \item \textbf{Use a model:} \emph{apply to map} makes a zoo model live (it registers on the tree
  and the map re-classifies); the base map card resets to the 4 classes. The zoo is usable, not just
  a catalogue.
  \item \textbf{Actions:} ``only for current view'' (the area filter), \emph{Show on map} for the
  extent, and \textbf{Publish} (single card or all).
\end{itemize}
\end{frame}

% ------------------------------------------------------------------
\begin{frame}{Describe it, prove it, map it}
A model card isn't just numbers; the user annotates it (the ``model cards in ML'' ask):
\begin{itemize}
  \item \textbf{Description, intended use, limitations:} what it's for and where it breaks.
  \item \textbf{Evidence:} how the classes were annotated (drone, field photos, \dots), the
  ``how do you know?'' the brief asked for.
  \item \textbf{Map each class to a standard:} an optional crosswalk to ESA WorldCover / USDA, picked
  from a built-in list (no codes to memorize, and any subset is fine), so a localized class stays
  interoperable without the model count blowing up.
\end{itemize}
\end{frame}

% ------------------------------------------------------------------
\begin{frame}{Quality feedback, and fixing skew}
Two guards the brief asked for, surfaced as live feedback on the cards:
\begin{itemize}
  \item \textbf{Spread (spatial diversity):} Shannon entropy of where a dataset's polygons fall,
  in $[0,1]$; flags ``100 polygons, all one district'' before it skews a model. (Diversity, not raw
  volume, is what generalizes on unseen regions.)
  \item \textbf{Class balance:} the support ratio across classes, flagged balanced / mild /
  imbalanced; if it's skewed, the user can retrain with \textbf{undersampling} or
  \textbf{oversampling} instead of the default class-weighting.
\end{itemize}
\end{frame}

% ------------------------------------------------------------------
\begin{frame}{Tea / non-tea: a separability check}
\begin{itemize}
  \item Tried the suggested \textbf{tea / non-tea} split of trees.
  \item Sampled Alpha Earth at 169 hand-marked polygons (105 tea / 64 non-tea), held out whole
  polygons, fit the same \texttt{LinearSVC} used everywhere else.
\end{itemize}
\vspace{0.3em}
\begin{center}\small
\begin{tabular}{@{}lr@{}}
\toprule
held-out accuracy & \textbf{0.934} \\
tea  & P 0.938 \quad R 0.961 \quad F 0.950 \\
non-tea & P 0.925 \quad R 0.884 \quad F 0.904 \\
\bottomrule
\end{tabular}
\end{center}
\end{frame}

% ------------------------------------------------------------------
\begin{frame}{What's in the catalogue}
\begin{columns}[T]
\begin{column}{0.5\textwidth}
\textbf{Models available}
\begin{itemize}\small
  \item Base map: greenery, water, built-up, barren (pan-India).
  \item Greenery split: crops, trees, shrubs.
  \item Barren split: barren, mining.
\end{itemize}
\end{column}
\begin{column}{0.5\textwidth}
\textbf{Datasets available}
\begin{itemize}\small
  \item Training: expert polygon sets (crops, trees, mining), the master expert pixel table,
  ESA WorldCover slices, extra water.
  \item Inference: the shared Alpha Earth feature source.
\end{itemize}
\end{column}
\end{columns}
\vspace{0.6em}
\begin{itemize}\footnotesize
  \item \textbf{11 cards} in all (3 model cards and 8 dataset cards), each \textbf{validated} against
  the schemas and pushed to the GitHub zoo.
  \item Every model card carries its training datasets, its inference feature source, and a bbox
  extent; held-out metrics live on each card.
\end{itemize}
\end{frame}

% ------------------------------------------------------------------
\begin{frame}
\centering
\Large Thank you
\end{frame}

\end{document}
